\documentclass[10pt,journal,compsoc]{IEEEtran}
\usepackage[utf8]{inputenc}
\usepackage{amsmath,amsfonts,amssymb}
\usepackage{graphicx}
\usepackage{cite}
\usepackage{algorithmic}
\usepackage{textcomp}

\begin{document}

\title{A Modular Generative Pipeline for High-Resolution Typographic Poster Synthesis}

\IEEEtitleabstractindextext{%
\begin{abstract}
The generation of professional, print-ready typographic posters using Artificial Intelligence presents unique challenges that cannot be solved via holistic, one-shot latent diffusion. This paper proposes a modular, multi-stage pipeline designed to decouple artistic background synthesis from rigid typographic rendering and print-density upscaling. We detail a sequential framework involving diffusion-based background generation, Multimodal Large Language Model (MLLM) guided layout planning, mask-guided typographic inpainting, and staged iterative upscaling. By separating semantic text rendering from textural generation, we establish a robust methodology for producing high-fidelity, 300 DPI posters free of upscaling hallucinations and illegible typography.
\end{abstract}
}

\maketitle
\IEEEdisplaynontitleabstractindextext
\IEEEpeerreviewmaketitle


\newtheorem{theorem}{Theorem}
\newtheorem{lemma}{Lemma}
\newtheorem{definition}{Definition}

\section{Introduction}
The commercial design industry requires posters to possess two distinct qualities: striking visual aesthetics and flawless, high-resolution typographic communication. Contemporary diffusion models (e.g., Stable Diffusion XL, FLUX) excel at generating complex aesthetics but fail critically when tasked with rendering coherent text and maintaining sharp structural edges at print resolutions ($\ge 300$ DPI).

The naive approach of "one-shot" prompting—attempting to generate the background, the subject, and the text simultaneously within a single latent pass—results in blended concepts, misspelled text, and chaotic layouts. To achieve professional quality, we propose abandoning the monolithic generation paradigm in favor of a strictly modular pipeline.

\section{The Modular Pipeline Architecture}
The proposed pipeline decomposes the poster synthesis task into four distinct stages, each handled by specialized computational nodes.

\subsection{Stage 1: Base Background Synthesis}
The first stage focuses entirely on the visual narrative. Using a high-parameter diffusion model, we synthesize a high-quality base image devoid of any typography.
\[ I_{base} = \mathcal{D}_{art}(P_{style}, P_{subject}, \Theta_{LoRA}) \]
where $\mathcal{D}_{art}$ is the diffusion operator, $P$ represents the textual prompts, and $\Theta_{LoRA}$ introduces Low-Rank Adaptations to enforce specific artistic aesthetics (e.g., minimalist, retro, cinematic).

\subsection{Stage 2: MLLM-Guided Layout and Typography Planning}
Once $I_{base}$ is generated, the spatial layout must be computed. We employ a Multimodal Large Language Model (MLLM) to analyze the negative space within $I_{base}$ and predict optimal bounding boxes $B = \{b_1, b_2, \dots, b_n\}$ for text placement. The MLLM enforces typographic hierarchy (Title, Subtitle, Copy) and prevents occlusion of critical visual subjects.

\subsection{Stage 3: Mask-Guided Typographic Inpainting}
Standard text overlays applied in post-processing often appear disconnected from the image lighting. To achieve compositional harmony, we utilize mask-guided inpainting.
Let $M_{text}$ be the binary mask derived from bounding boxes $B$. We employ a specialized text-rendering diffusion node (e.g., GlyphControl) to generate the text specifically within $M_{text}$:
\[ I_{comp} = \mathcal{D}_{text}(I_{base}, M_{text}, P_{text}) \]
This ensures the typography inherits the global lighting, shadows, and textural grain of the background, resulting in a cohesive integration.

\section{Iterative High-Density Upscaling}
The standard output resolution of latent diffusion models is typically $1024 \times 1024$ pixels, which is entirely insufficient for physical print media. A standard $24'' \times 36''$ poster at 300 DPI requires a resolution of $7200 \times 10800$ pixels.

\subsection{The Hallucination Problem in Upscaling}
Aggressive, single-pass upscaling (e.g., $8\times$ magnification) using generative upscalers introduces severe artifacts. Because generative upscalers (such as ControlNet Tile) attempt to "invent" detail where none exists, a single massive upscale often hallucinates unwanted micro-textures—turning smooth gradients into noise, or distorting precise typographic edges.

\subsection{Staged Tile-Based Upscaling Strategy}
To solve this, we formalize a staged, iterative upscaling methodology:
\begin{enumerate}
    \item \textbf{First Pass (Structural):} A conservative $2\times$ upscale using a vector-preserving or gentle upscaler (e.g., RealESRGAN) to establish crisp edges on the typography without altering the artistic texture.
    \item \textbf{Second Pass (Textural):} A subsequent tiled diffusion upscale with a low denoising strength ($\lambda_{denoise} < 0.35$). The image is divided into overlapping tiles $T_i$, processed individually, and blended using a Gaussian falloff to prevent visible seams.
\end{enumerate}
This iterative approach ensures that high-frequency details are gradually refined rather than abruptly hallucinated, preserving the integrity of both the typography and the underlying artwork.

\section{Print-Ready Standardization}
The final composite image $I_{final}$ must be prepared for physical reproduction. This requires a formal color-space transformation from the generative RGB spectrum to a subtractive CMYK profile, ensuring color accuracy when physical ink is applied to paper. Additionally, bleed margins must be computationally appended to the borders to account for mechanical trimming errors during the printing process.

\section{Conclusion}
This paper establishes a rigorous, modular pipeline for AI-driven poster generation. By isolating background synthesis from typographic inpainting, and by formalizing an iterative, staged upscaling protocol, we bypass the inherent limitations of monolithic diffusion models. This framework provides a scalable, reliable pathway for generating high-fidelity, print-ready commercial art at resolutions that meet and exceed professional industry standards.


\end{document}
